\runninghead{Appointed Texts and Locators}
\section{Appointed Texts and Study Locators}

\noindent\textbf{Standing notice on rights and study English.}
The approved proclaimed liturgical texts of the United States --- the \work{Roman Missal}, Third Edition (English translation copyright ICEL, 2010/2011), the protected modern Latin bodies of the composed Missal orations, Prefaces, and solemn blessings, and the \work{Lectionary for Mass for Use in the Dioceses of the United States of America} (New American Bible and Revised Grail / Abbey Psalms) --- are protected works and are neither reproduced nor closely paraphrased here. All English Scripture passages printed below come from the public-domain Douay--Rheims (Challoner) translation and serve strictly as \textbf{public-domain study English}, never as the approved proclaimed English of the liturgy. All Latin biblical texts below follow the public-domain Clementine Vulgate; the Latin Communion Antiphon is additionally verified against the secondary 2002 digital reproduction of the third typical edition of the \work{Missale Romanum}. Bracketed italics (\notread{like this}) show complete carrier verses where the Lectionary omits a subverse or clause, so the reader can see both the exact liturgical selection and its immediate canonical seam.

\medskip
\noindent\textbf{Vulgate and Hebrew numbering.}
Two numbering shifts affect this Sunday's Scripture texts. First, Psalms~25, 80, and~85 in Hebrew and U.S.\ Lectionary numbering correspond to Psalms~24, 79, and~84 in the Latin Vulgate and Douay--Rheims; within Psalm~80(79) and Psalm~85(84), individual verse numbers match one-to-one because both the Hebrew text and the Vulgate count the psalm superscription as verse~1. Second, in Isaiah, the plea to rend the heavens stands at Isaiah~63:19b in the Hebrew text and U.S.\ Lectionary, but begins chapter~64 as Isaiah~64:1 in the Vulgate and Douay--Rheims, shifting every subsequent verse of chapter~64 by one (so Lectionary Isa.~64:2b--7 equals Vulgate/Douay--Rheims Isa.~64:3b--8). The \hyperref[sec:date-location]{Scriptural Date and Location} appendix follows the Vulgate numbering of the inspected Latin liturgical controls; both systems are kept side by side throughout the commentary.

\subsection{Entrance Antiphon}\label{proper-entrance-antiphon}
\properlocus{Identifying Latin opening: \emph{Ad te levavi}. Cataloged biblical basis: Ps.~25:1--3 (Vg.~Ps.~24:1--3). Exact Missal citation cue and editorial truncation were not collated offline from the 2002/2008 Missal or 2010 ICEL Antiphonary; protected ICEL English is omitted. Under GIRM~48 (U.S.\ adaptation), four sung options are available, or the Missal antiphon is recited when there is no singing. Complete public-domain biblical carrier verses for Vg.~Ps.~24:1--3:}
\begin{latinproper}
\notread{In finem. Psalmus David.} Ad te, Domine, levavi animam meam: Deus meus, in te confido; non erubescam. Neque irrideant me inimici mei: etenim universi qui sustinent te, non confundentur.
\end{latinproper}
\begin{englishwitness}{Douay--Rheims (Challoner) study English --- Ps.~24:1--3 (Heb.~25:1--3)}
\notread{Unto the end, a psalm for David.} To thee, O Lord, have I lifted up my soul. In thee, O my God, I put my trust; let me not be ashamed. Neither let my enemies laugh at me: for none of them that wait on thee shall be confounded.
\end{englishwitness}

\subsection{Collect}\label{proper-collect}
\properlocus{Identifying Latin opening: \emph{Da, quaesumus, omnipotens Deus}. Missal-owned composed oration shared across Years~A, B, and~C (\work{Dominica I Adventus}, \work{Missale Romanum}, 2002/2008; U.S.\ English third edition, 2011, with the February 2021 Collect conclusion emendation). Protected modern Latin and ICEL English bodies are omitted.}
Because the modern Latin and English Collect is withheld under rights limits, no modern wording or close paraphrase is given here. For historical comparison only --- never as a substitute for the modern prayer --- the ancient \work{Sacramentarium Gelasianum Vetus} (Book~II, lxxxiv, \work{Postcommunio}~683, Wilson 1894, p.~218, inspected in an uncollated raw optical character recognition copy) preserves a patristic-era Roman antecedent with the same opening words. One word immediately before \emph{voluntatem} is unreadable in that raw copy and is left marked honestly rather than silently supplied:
\begin{latinproper}
Da, quaesumus, omnipotens Deus, cunctae familiae tuae \textnormal{[unreadable word in raw OCR]} voluntatem in Christo Filio tuo, Domino nostro venienti in operibus iustis aptos occurrere: et eius dexterae sociati regnum mereantur possidere caeleste. Per.
\end{latinproper}
In plain literal translation of the readable clauses, that ancient Roman prayer asks almighty God to grant his whole family a will to go forth fittingly in righteous deeds to meet Christ his Son our Lord at his coming, so that, joined to his right hand, they may be counted worthy to possess the heavenly kingdom. Exact lexical alignment with the modern 2002/2008 Missal was not collated offline.

\subsection{First Reading}\label{proper-first-reading}
\properlocus{OLM n.~2 / U.S.\ Lectionary No.~2: Isa.~63:16b--17, 19b; 64:2b--7 (Vulgate and Douay--Rheims: Isa.~63:16b--17; 64:1, 3b--8). Omitted: 63:16a, 63:18--19a (Vg.~63:18--19), and 64:1--2a (Vg.~64:2--3a). No shorter or longer form is appointed.}
\begin{latinproper}
\notread{Tu enim pater noster: et Abraham nescivit nos, et Israel ignoravit nos:} tu, Domine, pater noster, redemptor noster, a saeculo nomen tuum. Quare errare nos fecisti, Domine, de viis tuis; indurasti cor nostrum ne timeremus te? Convertere propter servos tuos, tribus haereditatis tuae. \notread{63:18--19 Vg.\ omitted} Utinam dirumperes caelos, et descenderes; a facie tua montes defluerent; \notread{64:2 Vg.\ omitted} \notread{Cum feceris mirabilia, non sustinebimus;} descendisti, et a facie tua montes defluxerunt. A saeculo non audierunt, neque auribus perceperunt; oculus non vidit, Deus, absque te, quae praeparasti exspectantibus te. Occurristi laetanti, et facienti justitiam; in viis tuis recordabuntur tui. Ecce tu iratus es, et peccavimus; in ipsis fuimus semper, et salvabimur. Et facti sumus ut immundus omnes nos, et quasi pannus menstruatae universae justitiae nostrae; et cecidimus quasi folium universi, et iniquitates nostrae quasi ventus abstulerunt nos. Non est qui invocet nomen tuum; qui consurgat, et teneat te. Abscondisti faciem tuam a nobis, et allisisti nos in manu iniquitatis nostrae. Et nunc, Domine, pater noster es tu, nos vero lutum; et fictor noster tu, et opera manuum tuarum omnes nos.
\end{latinproper}
\begin{englishwitness}{Douay--Rheims (Challoner) study English --- Isa.~63:16b--17; 64:1, 3b--8 (Heb.~63:16b--17, 19b; 64:2b--7)}
\notread{For thou art our father, and Abraham hath not known us, and Israel hath been ignorant of us:} thou, O Lord, art our father, our redeemer, from everlasting is thy name. Why hast thou made us to err, O Lord, from thy ways: why hast thou hardened our heart, that we should not fear thee? return for the sake of thy servants, the tribes of thy inheritance. \notread{63:18--19a Heb.\ omitted} O that thou wouldst rend the heavens, and wouldst come down: the mountains would melt away at thy presence. \notread{64:1 Heb.\ omitted} \notread{When thou shalt do wonderful things, we shall not bear them:} thou didst come down, and at thy presence the mountains melted away. From the beginning of the world they have not heard, nor perceived with the ears: the eye hath not seen, O God, besides thee, what things thou hast prepared for them that wait for thee. Thou hast met him that rejoiceth, and doth justice: in thy ways they shall remember thee: behold thou art angry, and we have sinned: in them we have been always, and we shall be saved. And we are all become as one unclean, and all our justices as the rag of a menstruous woman: and we have all fallen as a leaf, and our iniquities, like the wind, have taken us away. There is none that calleth upon thy name: that riseth up, and taketh hold of thee: thou hast hid thy face from us, and hast crushed us in the hand of our iniquity. And now, O Lord, thou art our father, and we are clay: and thou art our maker, and we all are the works of thy hands.
\end{englishwitness}

\subsection{Responsorial Psalm}\label{proper-responsorial-psalm}
\properlocus{OLM n.~2 / U.S.\ Lectionary No.~2: Ps.~80:2ac, 3b, 15--16, 18--19; Refrain: Ps.~80:4 (Vg.~Ps.~79:2ac, 3b, 15--16, 18--19; R.~v.~4). Omitted: vv.~2b, 3a, 5--14, and~17. Under GIRM~61, U.S.\ Lectionary No.~174 permits two seasonal Advent Common Psalms when sung --- Ps.~25(24) and Ps.~85(84) --- whose exact copyrighted U.S.\ refrain texts and verse spans remain uninspected offline.}
\begin{latinproper}
\textnormal{Refrain (v.~4):} Deus, converte nos, et ostende faciem tuam, et salvi erimus.\\
\textnormal{Stanza~1 (vv.~2ac, 3b):} Qui regis Israel, intende; \notread{qui deducis velut ovem Joseph.} Qui sedes super cherubim, manifestare \notread{coram Ephraim, Benjamin, et Manasse.} Excita potentiam tuam, et veni, ut salvos facias nos.\\
\textnormal{Stanza~2 (vv.~15--16):} Deus virtutum, convertere, respice de caelo, et vide, et visita vineam istam: et perfice eam quam plantavit dextera tua, et super filium hominis quem confirmasti tibi.\\
\textnormal{Stanza~3 (vv.~18--19):} Fiat manus tua super virum dexterae tuae, et super filium hominis quem confirmasti tibi. Et non discedimus a te: vivificabis nos, et nomen tuum invocabimus.
\end{latinproper}
\begin{englishwitness}{Douay--Rheims (Challoner) study English --- Ps.~79:2ac, 3b, 15--16, 18--19; R.~v.~4 (Heb.~Ps.~80)}
\textbf{Refrain (v.~4):} Convert us, O God: and shew us thy face, and we shall be saved.\\
\textbf{Stanza~1 (vv.~2ac, 3b):} Give ear, O thou that rulest Israel: \notread{thou that leadest Joseph like a sheep.} Thou that sittest upon the cherubims, shine forth \notread{Before Ephraim, Benjamin, and Manasses.} Stir up thy might, and come to save us.\\
\textbf{Stanza~2 (vv.~15--16):} Turn again, O God of hosts, look down from heaven, and see, and visit this vineyard: And perfect the same which thy right hand hath planted: and upon the son of man whom thou hast confirmed for thyself.\\
\textbf{Stanza~3 (vv.~18--19):} Let thy hand be upon the man of thy right hand: and upon the son of man whom thou hast confirmed for thyself. And we depart not from thee, thou shalt quicken us: and we will call upon thy name.
\end{englishwitness}

\subsection{Second Reading}\label{proper-second-reading}
\properlocus{OLM n.~2 / U.S.\ Lectionary No.~2: 1~Cor.~1:3--9. Continuous passage with no internal omissions and no shorter or longer form.}
\begin{latinproper}
Gratia vobis, et pax a Deo Patre nostro, et Domino Jesu Christo. Gratias ago Deo meo semper pro vobis in gratia Dei, quae data est vobis in Christo Jesu: quod in omnibus divites facti estis in illo, in omni verbo, et in omni scientia. Sicut testimonium Christi confirmatum est in vobis: ita ut nihil vobis desit in ulla gratia, exspectantibus revelationem Domini nostri Jesu Christi, qui et confirmabit vos usque in finem sine crimine, in die adventus Domini nostri Jesu Christi. Fidelis Deus: per quem vocati estis in societatem filii ejus Jesu Christi Domini nostri.
\end{latinproper}
\begin{englishwitness}{Douay--Rheims (Challoner) study English --- 1~Cor.~1:3--9}
Grace to you and peace, from God our father and from the Lord Jesus Christ. I give thanks to my God always for you, for the grace of God that is given you in Christ Jesus: That in all things you are made rich in him, in all utterance and in all knowledge; As the testimony of Christ was confirmed in you, So that nothing is wanting to you in any grace, waiting for the manifestation of our Lord Jesus Christ. Who also will confirm you unto the end without crime, in the day of the coming of our Lord Jesus Christ. God is faithful: by whom you are called unto the fellowship of his Son, Jesus Christ our Lord.
\end{englishwitness}

\subsection{Gospel Acclamation}\label{proper-gospel-acclamation}
\properlocus{Identifying Latin opening: \emph{Ostende nobis, Domine}. OLM n.~2 / U.S.\ Lectionary No.~2: Ps.~85:8 (Vg.~Ps.~84:8), framed by \emph{Alleluia} (GIRM~62).}
\begin{latinproper}
Ostende nobis, Domine, misericordiam tuam, et salutare tuum da nobis.
\end{latinproper}
\begin{englishwitness}{Douay--Rheims (Challoner) study English --- Ps.~84:8 (Heb.~85:8)}
Shew us, O Lord, thy mercy; and grant us thy salvation.
\end{englishwitness}

\subsection{Gospel}\label{proper-gospel}
\properlocus{OLM n.~2 / U.S.\ Lectionary No.~2: Mark~13:33--37. Continuous closing section of the Olivet Discourse, beginning immediately after the omitted verse Mark~13:32 (\emph{De die autem illo vel hora nemo scit, neque angeli in caelo, neque Filius, nisi Pater}: ``But of that day or hour no man knoweth, neither the angels in heaven, nor the Son, but the Father''). No internal omissions and no shorter or longer form.}
\begin{latinproper}
Videte, vigilate, et orate: nescitis enim quando tempus sit. Sicut homo qui peregre profectus reliquit domum suam, et dedit servis suis potestatem cujusque operis, et janitori praecepit ut vigilet, vigilate ergo (nescitis enim quando dominus domus veniat: sero, an media nocte, an galli cantu, an mane), ne, cum venerit repente, inveniat vos dormientes. Quod autem vobis dico, omnibus dico: Vigilate.
\end{latinproper}
\begin{englishwitness}{Douay--Rheims (Challoner) study English --- Mark~13:33--37}
Take ye heed, watch and pray. For ye know not when the time is. Even as a man who, going into a far country, left his house and gave authority to his servants over every work and commanded the porter to watch. Watch ye therefore (for you know not when the lord of the house cometh, at even, or at midnight, or at the cock crowing, or in the morning): Lest coming on a sudden, he find you sleeping. And what I say to you, I say to all: Watch.
\end{englishwitness}

\subsection{Prayer over the Offerings}\label{proper-prayer-over-the-offerings}
\properlocus{Identifying Latin opening: \emph{Suscipe, quaesumus, Domine}. Missal-owned composed oration shared across Years~A, B, and~C (\work{Dominica I Adventus}, \work{Missale Romanum}, 2002/2008; U.S.\ English third edition, 2011). Protected modern Latin and ICEL English bodies are omitted.}
While the modern Missal prayer itself is withheld, a verified comparison against the secondary 2002 digital reproduction of the \work{Missale Romanum} (pp.~82--83) confirms that the petition closely matches the ancient formula of the \work{Sacramentarium Veronense} (Section~XXVIIII, \work{Item alia}, marginal~372; Feltoe 1896, p.~74), differing in the petition body only by a single inflectional change (ancient \emph{temporalis} against secondary-2002 \emph{temporali}) plus editorial punctuation and the short ending cue (\emph{per.} against \emph{Per Christum.}); printed 2002 and printed 2008 correspondence remains unverified. The public-domain historical text from the \work{Veronense} reads:
\begin{latinproper}
Suscipe, quaesumus, Domine, munera quae de tuis offerimus collata beneficiis: et quod nostrae devotioni concedis effici temporalis, tuae nobis fiat praemium redemptionis aeternae. per.
\end{latinproper}
Literally rendered from the ancient Latin, the prayer asks the Lord to receive the gifts offered from his own bestowed benefits, and that what he grants to be accomplished for our temporal devotion may become for us the reward of his eternal redemption. No proper Offertory chant is printed in the Missal; under GIRM~74, an Offertory chant may accompany the rite under the norms for the Entrance chant.

\subsection{Seasonal Preface}\label{proper-preface}
\properlocus{Seasonal Preface of Advent (\work{Ordo Missae, Praefationes}, \work{Missale Romanum}, 2002/2008; GIRM~364). Protected modern Latin and ICEL English texts are omitted.}
Verified inspection of the \work{Ordo Missae} Preface inventory confirms that the third typical edition contains two Prefaces of Advent. Because the restricted Missal pages were not reopened offline, the exact Latin title, the calendar date-range rubric distinguishing the two Advent Prefaces, and the modern text appointed for the First Sunday of Advent remain uninspected offline and are neither quoted nor paraphrased. Under GIRM~365a--d, the seasonal Preface may be paired with Eucharistic Prayer~I (which carries no proper Advent insert in the \work{Communicantes} or \work{Hanc igitur}), Eucharistic Prayer~II, or Eucharistic Prayer~III; Eucharistic Prayer~IV is excluded because its Preface is invariable.

\subsection{Communion Antiphon}\label{proper-communion-antiphon}
\properlocus{Identifying Latin opening: \emph{Dominus dabit benignitatem}. Biblical basis: Ps.~85:13 (Vg.~Ps.~84:13). Verified public-domain Latin text collated against the secondary 2002 digital Missal reproduction (p.~83; correspondence with the printed 2008 edition remains unverified). Protected ICEL English is omitted; under GIRM~87 (U.S.\ adaptation), four sung options stand alongside recitation when there is no singing.}
\begin{latinproper}
Dominus dabit benignitatem, et terra nostra dabit fructum suum.
\end{latinproper}
\begin{englishwitness}{Douay--Rheims (Challoner) study English --- Ps.~84:13 (Heb.~85:13)}
\notread{For} the Lord will give goodness: and our earth shall yield her fruit.
\end{englishwitness}

\subsection{Prayer after Communion}\label{proper-prayer-after-communion}
\properlocus{Identifying Latin opening: \emph{Prosint nobis}. Missal-owned composed oration shared across Years~A, B, and~C (\work{Dominica I Adventus}, \work{Missale Romanum}, 2002/2008; U.S.\ English third edition, 2011). Protected modern Latin and ICEL English bodies are omitted.}
The held ancient sacramentary collections contain no direct ancestor for this prayer, and its modern Latin and English texts were not reopened offline. Commentary is bounded strictly to its Latin opening words (\emph{Prosint nobis}, ``May they profit us'') and its liturgical place closing the Communion Rite (\work{Ordo Missae} nn.~124--146).

\subsection{Concluding Blessing}\label{proper-concluding-blessing}
\properlocus{Ordinary blessing, or optional solemn blessing or prayer over the people (GIRM~90; \work{Ordo Missae} nn.~124--146; 2008 \work{Notitiae} nn.~503--504, \work{Ordo Missae} variations). Protected modern texts are omitted.}
Under GIRM~90, the Concluding Rites include the Priest's greeting and blessing, which on certain days and occasions may be enriched by a solemn blessing or a prayer over the people, followed by the dismissal (to which the 2008 emended reprint added three alternative formulas after \emph{Ite, missa est}, alongside the solemn-blessing invitation \emph{Inclinate vos ad benedictionem}, ``Bow down for the blessing''). The specific Advent solemn blessing numbering, Latin title, and protected wording were not inspected offline and are neither quoted nor paraphrased.
